\section{Conclusion}\label{sec:conclusion}

A* meets both requirements of Section~\ref{sec:requirements}. The cost so far
$g$ keeps it optimal, as in Dijkstra's algorithm, and the estimate $h$ directs
it toward the destination, as in greedy best-first search.

With a consistent heuristic, A* expands every vertex at most once and returns a
shortest path (Theorem~\ref{thm:consistent}). With an admissible but
inconsistent heuristic, a shortest path is guaranteed only when closed vertices
can be reopened (Theorem~\ref{thm:admissible}), and with an overestimating
heuristic no guarantee remains. The worst-case cost equals that of Dijkstra's
algorithm, but A* expands only vertices with $\delta(s,u) + h(u) \le C^{*}$
(Proposition~\ref{prop:expanded}). In our experiments it expanded 81 cells
against 137 on the grid and eight districts against 50 on the map of Bangladesh,
always returning a shortest path.

Since a larger admissible heuristic leaves fewer vertices to expand, choosing a
heuristic that is cheap and as close as possible to the true cost is the main
design decision when A* is applied to a new problem.

\paragraph{Code availability.} The programs that produce every trace, table and
experimental result of this report are available at \url{\repourl}.
